\section{Related Work}
\label{sec:related}

We position this paper at the intersection of five related but
distinct strands of work.

\subsection{Compact Recurrent Cells}

The standard LSTM~\cite{hochreiter1997lstm} and
GRU~\cite{cho2014gru} cells encode long-range temporal structure
effectively but are too heavy for kilobyte-class MCUs: a
modestly-sized LSTM ($H{=}64$, $d{=}3$) already exceeds
\SI{50}{\kilo\byte} of weights at FP32 precision, well above the
\SI{16}{\kilo\byte} Flash budget of our target device.

\fastgrnn{}~\cite{kusupati2018fastgrnn} addresses this gap with a
two-scalar gated cell ($\zeta, \nu$) backed by low-rank weight
matrices, achieving accuracy comparable to LSTM on sequence tasks
at one to two orders of magnitude fewer parameters. Complementary
non-recurrent EdgeML approaches include
Bonsai~\cite{kumar2017bonsai}, a shallow non-linear tree, and
ProtoNN~\cite{gupta2017protonn}, a prototype-based classifier.
The three together form the Microsoft Research EdgeML
toolkit~\cite{edgeml}. \fastgrnn{} is the only one of these three
that natively handles temporal input, which is essential for HAR.

\subsection{Quantization}

Post-training quantization (PTQ) compresses neural-network weights
and activations from FP32 to fixed-point or low-bit integer
formats. Jacob \emph{et al.}~\cite{jacob2018integer} introduce
integer-arithmetic-only inference with a per-tensor scale
formulation; Han \emph{et al.}~\cite{han2016deepcompression}
combine pruning, quantization, and Huffman coding to compress deep
networks by an order of magnitude.
Quantization-aware training~(QAT)~\cite{hubara2017qnn} avoids the
PTQ accuracy cliff at the price of full retraining. We use
per-tensor Q15 PTQ for both weights and -- with explicit activation
calibration -- intermediate tensors; the calibration step turns
out to be the dividing line between lossless and catastrophic
deployment (Section~\ref{sec:results:quant}), and is what allows
us to avoid the QAT machinery entirely.

\subsection{Sparsity and Pruning}

Iterative hard thresholding (IHT)~\cite{blumensath2009iht} is the
sparsification engine used by \fastgrnn{}: at each step the
top-$k$ magnitude entries of every weight tensor are retained and
the rest are zeroed, then the network is re-trained with the mask
fixed. Han \emph{et al.}~\cite{han2015pruning} demonstrate that
aggressive pruning preserves accuracy in larger
networks; Frankle and Carbin's lottery ticket
hypothesis~\cite{frankle2019lt} provides theoretical grounding for
why such sparse sub-networks remain trainable.

\subsection{tinyML on More Capable Targets}

The tinyML community~\cite{banbury2021mlperf} has converged on
ARM Cortex-M class targets with hundreds of kilobytes of SRAM,
hardware FPUs, and single-cycle multipliers.
MCUNet~\cite{lin2020mcunet} jointly designs a neural architecture
and an interpreter to maximize Cortex-M throughput.
CMix-NN~\cite{capotondi2020cmixnn} provides mixed-precision
kernels for memory-constrained edge devices.
CMSIS-NN~\cite{lai2018cmsisnn} from ARM accelerates common layers
via Cortex-M SIMD instructions, and TensorFlow Lite
Micro~\cite{david2021tflitemicro} offers a portable runtime layer.

None of these target the \msp{}-class device addressed in
this paper, which is one to two orders of magnitude smaller and
lacks a hardware multiplier entirely. To our knowledge, the present
work is the first end-to-end public reproduction of a gated
recurrent cell on an MCU with no hardware multiplier.

\subsection{Human Activity Recognition}

The UCI HAR~\cite{anguita2013public} and its transition-aware
extension \hapt{}~\cite{reyes2015transition} are the de facto
public benchmarks for accelerometer-based HAR.
DeepConvLSTM~\cite{ordonez2016deepconvlstm} established a strong
deep-learning baseline using stacked CNN and LSTM layers; broader
surveys by Wang \emph{et al.}~\cite{wang2019harsurvey} and
Demrozi \emph{et al.}~\cite{demrozi2020harsurvey} catalogue the
methodological landscape, both of which highlight the dynamic
\textsc{downstairs} class as a persistent failure
mode -- a finding we corroborate in
Section~\ref{sec:results:perclass}.
